% GENERATED FILE. Edit canonical lessons, then run npm run generate:books.
% canonical-source-hash: fnv1a64:b686821807e7d298
% canonical-lessons: KA-C149-ili2, KA-C149-seru, KA-C149-todu, KA-C149-kare, KA-C149-enisu

\chapter{To get down, To join, To wear, To call, To count}
\label{ch:ka-to-get-down-to-join-to-wear-to-call-to-count}
\hlchaptermodality{149}

\begin{chapteropening}
\textbf{By the end of this chapter:} \emph{I can say to get down, to join, to wear, to call and to count.}

Complete the last lesson of chapter 149: i can say to get down, to join, to wear, to call and to count.
\end{chapteropening}

\section[iḷi]{\kn{ಇಳಿ} (iḷi) — to get down, to get off}

\label{lesson:KA-C149-ili2}



\begin{quote}
Before the new one: say the Kannada for to teach, then the Kannada for to jump.
\end{quote}

\subsection*{You'll want to know: \kn{ಇಳಿ}}
\textbf{\kn{ಇಳಿ}} — \emph{iḷi} — \textquotedblleft{}to get down, to get off\textquotedblright{}.

\begin{grammarlens}[title={What you've built}]
One more word for this chapter.
\end{grammarlens}

\subsection*{Your turn}
\begin{itemize}
\raggedright
  \item \emph{Say it:} \emph{iḷi}
  \item \emph{Say it:} \emph{iḷi}, once more
  \item \emph{Say it:} say \emph{iḷi}
  \item \emph{Recall:} say the Kannada for to teach, then the Kannada for to jump, then say \emph{iḷi} again
\end{itemize}

\begin{tcolorbox}[breakable,colback=teal!4,colframe=teal!35!black,title={Before you move on}]
What does \kn{ಇಳಿ} mean? (\textquotedblleft{}To get down, to get off\textquotedblright{}.) Say it once more.
\end{tcolorbox}
\section[sēru]{\kn{ಸೇರು} (sēru) — to join, to reach}

\label{lesson:KA-C149-seru}



\begin{quote}
Before the new one: say the Kannada for to jump, then the Kannada for to get down.
\end{quote}

\subsection*{You'll want to know: \kn{ಸೇರು}}
\textbf{\kn{ಸೇರು}} — \emph{sēru} — \textquotedblleft{}to join, to reach\textquotedblright{}.

\begin{grammarlens}[title={What you've built}]
One more word for this chapter.
\end{grammarlens}

\subsection*{Your turn}
\begin{itemize}
\raggedright
  \item \emph{Say it:} \emph{sēru}
  \item \emph{Say it:} \emph{sēru}, once more
  \item \emph{Say it:} say \emph{sēru}
  \item \emph{Recall:} say the Kannada for to jump, then the Kannada for to get down, then say \emph{sēru} again
\end{itemize}

\begin{tcolorbox}[breakable,colback=teal!4,colframe=teal!35!black,title={Before you move on}]
What does \kn{ಸೇರು} mean? (\textquotedblleft{}To join, to reach\textquotedblright{}.) Say it once more.
\end{tcolorbox}
\section[toḍu]{\kn{ತೊಡು} (toḍu) — to wear}

\label{lesson:KA-C149-todu}



\begin{quote}
Before the new one: say the Kannada for to get down, then the Kannada for to join.
\end{quote}

\subsection*{You'll want to know: \kn{ತೊಡು}}
\textbf{\kn{ತೊಡು}} — \emph{toḍu} — \textquotedblleft{}to wear\textquotedblright{}.

\begin{grammarlens}[title={What you've built}]
One more word for this chapter.
\end{grammarlens}

\subsection*{Your turn}
\begin{itemize}
\raggedright
  \item \emph{Say it:} \emph{toḍu}
  \item \emph{Say it:} \emph{toḍu}, once more
  \item \emph{Say it:} say \emph{toḍu}
  \item \emph{Recall:} say the Kannada for to get down, then the Kannada for to join, then say \emph{toḍu} again
\end{itemize}

\begin{tcolorbox}[breakable,colback=teal!4,colframe=teal!35!black,title={Before you move on}]
What does \kn{ತೊಡು} mean? (\textquotedblleft{}To wear\textquotedblright{}.) Say it once more.
\end{tcolorbox}
\section[kare]{\kn{ಕರೆ} (kare) — to call}

\label{lesson:KA-C149-kare}



\begin{quote}
Before the new one: say the Kannada for to join, then the Kannada for to wear.
\end{quote}

\subsection*{You'll want to know: \kn{ಕರೆ}}
\textbf{\kn{ಕರೆ}} — \emph{kare} — \textquotedblleft{}to call\textquotedblright{}.

\begin{grammarlens}[title={What you've built}]
One more word for this chapter.
\end{grammarlens}

\subsection*{Your turn}
\begin{itemize}
\raggedright
  \item \emph{Say it:} \emph{kare}
  \item \emph{Say it:} \emph{kare}, once more
  \item \emph{Say it:} say \emph{kare}
  \item \emph{Recall:} say the Kannada for to join, then the Kannada for to wear, then say \emph{kare} again
\end{itemize}

\begin{tcolorbox}[breakable,colback=teal!4,colframe=teal!35!black,title={Before you move on}]
What does \kn{ಕರೆ} mean? (\textquotedblleft{}To call\textquotedblright{}.) Say it once more.
\end{tcolorbox}
\section[eṇisu]{\kn{ಎಣಿಸು} (eṇisu) — to count}

\label{lesson:KA-C149-enisu}



\begin{quote}
Before the new one: say the Kannada for to wear, then the Kannada for to call.

Say the ordinals past eight: \textbf{\kn{ಒಂಬತ್ತನೆಯ}}, \textbf{\kn{ಹತ್ತನೆಯ}}, \textbf{\kn{ಹನ್ನೊಂದನೆಯ}}. Then the spoken \textbf{\kn{ಮೊದಲನೇ}}, and read the sign form \textbf{\kn{೧ನೇ}}.
\end{quote}

\subsection*{You'll want to know: \kn{ಎಣಿಸು}}
\textbf{\kn{ಎಣಿಸು}} — \emph{eṇisu} — \textquotedblleft{}to count\textquotedblright{}.

\begin{grammarlens}[title={What you've built}]
One more word for this chapter.
\end{grammarlens}

\subsection*{Your turn}
\begin{itemize}
\raggedright
  \item \emph{Say it:} \emph{eṇisu}
  \item \emph{Say it:} \emph{eṇisu}, once more
  \item \emph{Say it:} say \emph{eṇisu}
  \item \emph{Recall:} say the Kannada for to wear, then the Kannada for to call, then say \emph{eṇisu} again
\end{itemize}

\begin{tcolorbox}[breakable,colback=teal!4,colframe=teal!35!black,title={Before you move on}]
What does \kn{ಎಣಿಸು} mean? (\textquotedblleft{}To count\textquotedblright{}.) Say it once more.
\end{tcolorbox}
